{
\mbox{}
\thispagestyle{empty}

\vfill

\footnotesize
\noindent{}\emph{Pegue o primeiro texto que encontrar: \emph{texto-memória},
\emph{texto-qualquer}, \emph{texto-nenhum}. Ele sempre será um mapa. Foi
o que me ocorreu ao ler o retrato de adeus e de amor de minha trisavó
Dvora. % \footnote{Dvora é mãe de Aba, que é pai de Sarah Dora, que é mãe de Luciana, que é mãe de Tamara --- que nasceu na cidade de São Paulo, em 1999.}. 
  Eu o acolhi em mim como uma espécie de guia, uma casa, um
mundo por completo. Agora é hora de pousar nas páginas, dentro delas. Trazer as palavras para perto, fazê-las falar o que queremos dizer. O que escutamos ao colar o ouvido no espelho?}

\pagebreak
}

\chapter*{Introdução\smallskip\subtitulo{Lorem Ipsum Dolor Sit Amet}}
\addcontentsline{toc}{chapter}{Introdução, \textit{por Tamara Crespin}}

\begin{flushright}
\textsc{tamara crespin}
\end{flushright}

\noindent{}Ainda prefiro aquelas coisas sobre as quais não se sabe ao certo o que foi inventado. Diz-se que existem apenas sete histórias possíveis de serem contadas no mundo, e que o que muda entre elas é a capacidade de rearranjo de cada coletor, de cada recorte, de cada ritmo, de cada maneira de entrelaçar essas narrativas. Não saberia dizer quais são essas sete histórias. Sei, porém, que continuamos escrevendo, de diferentes maneiras, sempre o mesmo livro.

Em \emph{A ensaificação de tudo}, a escritora norte-americana Christy Wampole afirma que ``O ensaísta está interessado em pensar sobre si pensando sobre as coisas''.\footnote{\textsc{wampole}, \emph{A ensaificação de tudo}.} Ao fim e ao cabo, esta pesquisa é também uma espécie de ensaio que, assim como a própria diáspora, sustenta-se como tentativa de solução, ainda que na ausência de soluções. São impulsos e propostas para ``pensar sobre si pensando sobre as coisas''\footnote{\emph{Ibidem}.} e, a partir desse movimento, construir outras territorialidades e frequências comunitárias. Assim sendo,

\begin{quote}
Esse não é um livro de história. A escolha que nele se encontrará não
seguiu outra regra mais importante do que meu gosto, meu prazer, uma
emoção, o riso, a surpresa, um certo assombro ou qualquer outro
sentimento, do qual teria dificuldades, talvez, em justificar a
intensidade, agora que o primeiro momento da descoberta
passou\footnote{\textsc{foucault}, \emph{A vida dos homens infames}, p.\,203.}.
\end{quote}

Busca-se aqui compreender que as \emph{perspectivas judaico-diaspóricas sobre o lugar} estão entremeadas de signos, símbolos, línguas e arquiteturas em suas mais variadas escalas, nas quais se permitem o nascimento e a consolidação de identidades fluidas, atravessadas pela indeterminabilidade e, por isso, passíveis de apropriação e de pertencimento.

São as fagulhas, as memórias-palavras que perduram, que vagueiam pelo ar e são levadas pelo velho sopro do vento nômade da diáspora. Ao pousarem, criam territórios que, mesmo temporários, possibilitam a construção de lugares próprios de sentido e identificação. Mais do que espaços físicos, são lugares que se manifestam por meio de sentidos identitários e comunitários, através das línguas, dos objetos culturais e dos ritos religiosos. São espaços de trânsito, do indeterminado, do fluido e do híbrido, que atravessam, em suas mais diversas formas, as noções espaço-temporais judaicas.

Essas geografias de identidades serão investigadas no decorrer deste texto por meio dos conceitos judaicos de:

\begin{itemize}
\item \textsc{bs'd}; 
\item Kipá; 
\item Mezuzá; 
\item Eruv; 
\item Makom. 
\end{itemize}

Transversalmente, procura-se compreender a página, as palavras e as línguas diaspóricas\footnote{As línguas também possuem biografia. O ladino conserva memórias de infância de um espanhol medieval. As pessoas podem esquecer, mas as línguas não. Elas mantêm as alegrias, as dores e as feridas de tudo o que há passado.} como elementos essenciais à criação de corpos identitários e comunitários; a \emph{Shechiná} --- lugar de habitação, crença e presença divina --- como uma territorialidade temporária; e a diáspora como limite e limiar em si mesma, um espaço eternamente transitório.

Rabi Tarfon\footnote{Rabi Tarfon viveu entre os séculos \textsc{i}--\textsc{ii} d.E.\,C. e foi membro da terceira geração de
  estudiosos da \emph{Mishná}.}
costumava dizer: ``Não lhe é exigido que complete a tarefa, mas
você não é livre para escapar dela''\footnote{\emph{Pirkei Avot} 2:16.}. Essa constante elaboração é o que
impossibilita o esgotamento de sentidos sobre o ser judeu. Por fim, já
que estaremos sempre a escrever o mesmo livro, essa investigação
sustenta-se enquanto um ensaio identitário, sem a pretensão de alicerçar
e fixar imagens judaicas.

Rabi Tarfon\footnote{Rabi Tarfon foi membro da terceira geração de estudiosos da \emph{Mishná}, viveu entre os séculos \textsc{i} e \textsc{ii} d.E.C.} costumava dizer: ``Não lhe é exigido que complete a tarefa, mas você não é livre para escapar dela''.\footnote{\emph{Pirkei Avot} 2:16.} É justamente essa tarefa, sempre inacabada, que impossibilita o esgotamento dos sentidos do ``ser judeu''. % Por fim, já que estaremos sempre a escrever o mesmo livro, esta investigação se sustenta como um ensaio identitário, sem a pretensão de alicerçar ou fixar imagens judaicas.
Esta breve, dispersa e pequena introdução é apenas uma abertura de
portas. \emph{I ke io lo mire}.\footnote{Do ladino, ``e que eu o veja''.}